% BuzzerModule — ブザー出力モジュール
\subsection{BuzzerModule}

\begin{apibox}{BuzzerModule --- ブザー出力モジュール}
\textbf{定義ファイル:} \texttt{lib/BuzzerModule/BuzzerModule.h / .cpp}\\
\textbf{分類:} 出力モジュール（\texttt{updateOutput()}のみオーバーライド）\\
\textbf{対象デバイス:} パッシブブザー（PWM駆動）

LEDCによるPWMトーン再生を行うモジュール。
周波数と再生時間を指定でき、時間経過後に自動停止する。
\end{apibox}

\subsubsection{機能概要}

\begin{itemize}
    \item LEDC PWMによるトーン生成（任意の周波数）
    \item 再生時間指定（\texttt{durationMs}）による自動停止
    \item 再生リクエストフラグ（\texttt{requestPlay}）による制御
    \item \texttt{deinit()}によるPWMチャネル解放
\end{itemize}

% --- Config ---
\subsubsection{Config構造体}

\begin{funcblock}{BuzzerConfig}
ブザーの接続ピンとPWMチャネルを定義する。

\begin{longtable}{l l l p{5.0cm}}
\toprule
\textbf{メンバ} & \textbf{型} & \textbf{制約・既定値} & \textbf{説明} \\
\midrule
\endhead
\texttt{pin}        & \texttt{uint8\_t} & ---         & ブザー出力GPIOピン番号 \\
\texttt{pwmChannel} & \texttt{uint8\_t} & 0--15       & LEDCチャネル番号。他モジュールと重複しないこと \\
\bottomrule
\end{longtable}
\end{funcblock}

% --- Data ---
\subsubsection{Data構造体}

\begin{funcblock}{BuzzerData}
ブザーの再生パラメータを保持する。ロジックフェーズで設定し、出力フェーズで消費される。

\begin{longtable}{l l l p{4.5cm}}
\toprule
\textbf{メンバ} & \textbf{型} & \textbf{初期値} & \textbf{説明} \\
\midrule
\endhead
\texttt{frequency}   & \texttt{uint16\_t}      & \texttt{0}     & 再生周波数 [Hz]。\texttt{0}で消音 \\
\texttt{durationMs}  & \texttt{unsigned long}  & \texttt{0}     & 再生時間 [ms]。\texttt{0}で無制限（停止は\texttt{frequency = 0}で行う） \\
\texttt{requestPlay} & \texttt{bool}           & \texttt{false} & 再生リクエストフラグ。ロジックフェーズで\texttt{true}にセットする \\
\bottomrule
\end{longtable}
\end{funcblock}

% --- メソッド ---
\subsubsection{メソッド}

\begin{funcblock}{BuzzerModule(const BuzzerConfig\& config)}
\begin{tabularx}{\textwidth}{l X}
\toprule
\textbf{項目} & \textbf{内容} \\
\midrule
引数   & \texttt{config} --- \texttt{BuzzerConfig}へのconst参照 \\
動作   & Configを\texttt{\_config}にコピー保持する \\
\bottomrule
\end{tabularx}
\end{funcblock}

\begin{funcblock}{bool init()}
\begin{tabularx}{\textwidth}{l X}
\toprule
\textbf{項目} & \textbf{内容} \\
\midrule
戻り値   & \texttt{bool} --- 常に\texttt{true} \\
動作     & LEDCチャネルの初期設定を行い、ブザーを消音状態にする \\
\bottomrule
\end{tabularx}
\end{funcblock}

\begin{funcblock}{void updateOutput(SystemData\& data)}
\begin{tabularx}{\textwidth}{l X}
\toprule
\textbf{項目} & \textbf{内容} \\
\midrule
読み取り元     & \texttt{data.buzzer} \\
タイミング制御 & 毎ループ実行 \\
動作           & \texttt{requestPlay}が\texttt{true}の場合、指定周波数でPWMトーンを開始し、\texttt{durationMs > 0}なら時間計測を開始。\texttt{requestPlay}を\texttt{false}にクリア。再生中に\texttt{durationMs}経過で自動停止 \\
\bottomrule
\end{tabularx}
\end{funcblock}

\begin{funcblock}{void deinit()}
\begin{tabularx}{\textwidth}{l X}
\toprule
\textbf{項目} & \textbf{内容} \\
\midrule
動作   & PWMトーンを停止し、LEDCチャネルをデタッチする \\
\bottomrule
\end{tabularx}
\end{funcblock}

\begin{notebox}[\texttt{requestPlay}のライフサイクル]
\texttt{requestPlay}はロジックフェーズで\texttt{true}にセットし、
\texttt{updateOutput()}内で消費（\texttt{false}にクリア）される。
毎ループ\texttt{true}にし続けると再生が毎ループリスタートされるため、
単発再生の場合は1回だけセットすること。
\end{notebox}

% --- 使用例 ---
\subsubsection{使用例}

\begin{lstlisting}[style=cppstyle, caption={BuzzerModuleの使用例}]
// ProjectConfig.h
const BuzzerConfig BUZZER_CONFIG = {
    .pin        = 47,  // GPIO47: ブザー出力
    .pwmChannel = 6,   // LEDCチャネル6
};

// ロジックフェーズでの使用
void applyPattern(SystemData& data) {
    // ボタン押下で「ピッ」と鳴らす
    if (data.button.justPressed) {
        data.buzzer.frequency   = 1000;  // 1kHz
        data.buzzer.durationMs  = 100;   // 100ms
        data.buzzer.requestPlay = true;
    }
}
\end{lstlisting}
